\subsection{Modes of Convergence}

% **Role:** compare pointwise, almost-everywhere, uniform, almost-uniform, in-measure, and $L^p$ convergence.

Throughout this subsection $(X,\mathcal{M},\mu)$ is a fixed measure space and, unless otherwise stated, functions are real-valued and measurable. The point is not to rank convergence modes linearly, but to know which hypotheses convert one mode into another.

\begin{definition} \label{def:modes-pointwise-uniform-ae} \hlt{Pointwise, Uniform, and Almost-Everywhere Convergence}\\
Let $(f_n)$ and $f$ be functions on $X$.
\begin{enumerate}[label=(\roman*)]
\setlength{\itemsep}{0pt}
\item $f_n\to f$ \hlt{pointwise} if $\forall x\in X$, $\varepsilon>0$, $\exists N(\varepsilon,x)$ s.t. $n\geq N(\varepsilon,x)$ implies $\abs{f_n(x)-f(x)}<\varepsilon$.
\item $f_n\to f$ \hlt{uniformly} if $\forall \varepsilon>0$ $\exists N(\varepsilon)$ s.t. $n\geq N(\varepsilon)$ implies $\abs{f_n(x)-f(x)}<\varepsilon$ $\forall x\in X$.
\item $f_n\to f$ \hlt{almost everywhere} if $\exists A\in\mathcal{M}$ with $\mu(A)=0$ s.t. $\forall x\in X\setminus A$, $\varepsilon>0$,\\
$\exists N(\varepsilon,x)$ s.t. $n\geq N(\varepsilon,x)$ implies $\abs{f_n(x)-f(x)}<\varepsilon$.
\end{enumerate}
\end{definition}

\begin{proposition} \label{prop:basic-mode-implications} \hlt{First Implications and First Failures}\\
Let $f_n,f:X\to\R$.
\begin{enumerate}[label=(\roman*)]
\setlength{\itemsep}{0pt}
\item $f_n\rightrightarrows f\Rightarrow f_n\to f$ pointwise $\Rightarrow f_n\to f$ a.e.; converses fail.
\item $\abs{X}<+\infty$ and $f_n\to f$ pointwise $\Rightarrow f_n\rightrightarrows f$.
\item $\mu(A)=0\Rightarrow A=\emptyset$ and $f_n\to f$ a.e. $\Rightarrow f_n\to f$ pointwise.
\end{enumerate}
\begin{proof}
\begin{enumerate}[label=(\roman*)]
\setlength{\itemsep}{0pt}
\item If $f_n\rightrightarrows f$, then $\forall\varepsilon>0\ \exists N_\varepsilon$ such that $\forall n\geq N_\varepsilon\ \forall x\in X$, $\abs{f_n(x)-f(x)}<\varepsilon$; hence pointwise convergence with $N(\varepsilon,x)=N_\varepsilon$. Pointwise convergence is a.e. convergence with $A=\emptyset$.\\
Conversely, on $([0,1],\lambda)$, $x^n\to\mathbf{1}_{\{1\}}$ pointwise but not uniformly, and $\mathbf{1}_{\{0\}}\to0$ a.e. but not pointwise.
\item If $X$ is finite and $f_n\to f$ pointwise, then $\forall\varepsilon>0\ \forall x\in X$, choose $N(\varepsilon,x)$.\\
Set $N_\varepsilon=\max\limits_{x\in X}N(\varepsilon,x)$; then $\forall n\geq N_\varepsilon\ \forall x\in X$, $\abs{f_n(x)-f(x)}<\varepsilon$, so $f_n\rightrightarrows f$.
\item If $f_n\to f$ a.e., then $\exists A\in\mathcal{M}$ with $\mu(A)=0$ and $\forall x\in X\setminus A$, $f_n(x)\to f(x)$.\\
Since every null set is empty, $A=\emptyset$, so $f_n\to f$ pointwise.
\end{enumerate}
\hlt{Structural Commentary}: Uniformity is bounded $x$-dependence of $N(\varepsilon,x)$.\\
\hlt{Key Insight}: Uniformity controls indices; a.e. convergence controls null exceptions.
\end{proof}
\end{proposition}

\begin{proposition} \label{prop:uniform-to-lp-on-finite-measure} \hlt{Uniform Convergence Gives $L^p$ Convergence on Finite Measure Spaces}\\
Let $1\leq p<\infty$, $\mu(X)<+\infty$, and $f_n\in L^p(\mu)$. If $f_n\rightrightarrows f$ on $X$, then $f\in L^p(\mu)$ and $\norm{f_n-f}_p\to0$.
\begin{proof}
For $\varepsilon>0$, choose $N$ with $n\geq N\Rightarrow \abs{f_n-f}\leq\varepsilon$ on $X$.\\
Then $n\geq N\Rightarrow \norm{f_n-f}_p\leq(\int_X\varepsilon^p\,d\mu)^{\frac{1}{p}}=\varepsilon\mu(X)^{\frac{1}{p}}$, so $\norm{f_n-f}_p\to0$.\\
Taking $\varepsilon=1$, choose $n_0$ with $\abs{f_{n_0}-f}\leq1$ on $X$; hence $f-f_{n_0}\in L^p(\mu)$ and $f=f_{n_0}+(f-f_{n_0})\in L^p(\mu)$.\\
\hlt{Structural Commentary}: Uniform error becomes $L^p$ error after multiplication by $\mu(X)^{\frac{1}{p}}$.\\
\hlt{Key Insight}: Finite measure converts a pointwise sup-bound into an integral bound.
\end{proof}
\end{proposition}

\begin{theorem} \label{thm:dominated-ae-to-lp} \hlt{Dominated Almost-Everywhere Convergence Gives $L^p$ Convergence}\\
Let $1\leq p<\infty$, $f_n\in L^p(\mu)$, and $f_n\to f$ a.e. Suppose $\exists g\in L^p(\mu)$ such that $\forall n,\ \abs{f_n}\leq g$ a.e.
\begin{enumerate}[label=(\roman*)]
\setlength{\itemsep}{0pt}
\item $f\in L^p(\mu)$.
\item $\norm{f_n-f}_p\to0$.
\end{enumerate}
\begin{proof}
\begin{enumerate}[label=(\roman*)]
\setlength{\itemsep}{0pt}
\item Choose $N\in\mathcal M$ with $\mu(N)=0$ such that on $X\setminus N$, $f_n\to f$ and $\abs{f_n}\leq g$ $\forall n$.\\
Then $\abs f\leq g$ a.e., so $\abs f^p\leq g^p\in\mathscr{L}^1(\mu)$.
\item Since $\abs{f_n-f}^p\to0$ a.e. and $\abs{f_n-f}^p\leq(2g)^p\in\mathscr{L}^1(\mu)$ a.e.,\\
Theorem \ref{thm:dominated-convergence-lebesgue} yields $\norm{f_n-f}_p^p=\int_X\abs{f_n-f}^p\,d\mu\to0$.
\end{enumerate}
\hlt{Structural Commentary}: Apply ordinary dominated convergence to the $p$th-power error.\\
\hlt{Key Insight}: A single $L^p$ envelope controls both membership of the limit and convergence in mean.
\end{proof}
\end{theorem}

\begin{corollary} \label{cor:bounded-ae-finite-measure-lp} \hlt{Uniform Boundedness on Finite Measure Spaces}\\
Let $\mu(X)<+\infty$, $0\leq K<+\infty$, and $f_n\to f$ a.e. Assume $\forall n,\ \abs{f_n}\leq K$ a.e. Then for every $1\leq p<\infty$:
\begin{enumerate}[label=(\roman*)]
\setlength{\itemsep}{0pt}
\item $f\in L^p(\mu)$.
\item $\norm{f_n-f}_p\to0$.
\end{enumerate}
\begin{proof}
Fix $1\leq p<\infty$ and set $g\defeq K\mathbf{1}_X$. Then $\int_Xg^p\,d\mu=K^p\mu(X)<+\infty$, so $g\in L^p(\mu)$ and $\abs{f_n}\leq g$ a.e.
\begin{enumerate}[label=(\roman*)]
\setlength{\itemsep}{0pt}
\item Since $f_n\to f$ a.e. and $\abs{f_n}\leq g\in L^p(\mu)$ a.e., Theorem \ref{thm:dominated-ae-to-lp}(i) gives $f\in L^p(\mu)$.
\item The same hypotheses and Theorem \ref{thm:dominated-ae-to-lp}(ii) give $\norm{f_n-f}_p\to0$.
\end{enumerate}
\hlt{Structural Commentary}: Finite measure turns a uniform a.e. bound into an $L^p$ envelope.\\
\hlt{Key Insight}: Bounded a.e. convergence on finite measure is dominated convergence.
\end{proof}
\end{corollary}

\begin{example} \label{ex:modes-uniform-not-lp-infinite-measure} \hlt{Uniform Need Not Imply $L^p$ on Infinite Measure}\\
On $(\R,\mathcal{B}(\R),\lambda)$, fix $1\leq p<\infty$ and set $f_n\defeq n^{-\frac{1}{p}}\mathbf{1}_{[0,n]}$. Then $\norm{f_n}_\infty=n^{-\frac{1}{p}}\to0$, so $f_n\rightrightarrows0$.\\
However $\norm{f_n}_p^p=\int_0^n n^{-1}\,d\lambda=1$, hence $\norm{f_n}_p\nrightarrow0$.\\
\hlt{Structural Commentary}: The support grows exactly enough to preserve $L^p$ mass.\\
\hlt{Key Insight}: Uniform smallness alone does not control integrals on infinite measure spaces.
\end{example}

\begin{example} \label{ex:modes-ae-measure-not-lp-finite-measure} \hlt{Almost Everywhere and In-Measure Need Not Imply $L^p$}\\
On $[0,1]$ with Lebesgue measure, set $f_n\defeq n\mathbf{1}_{[\frac{1}{n},\,\frac{2}{n}]}$.\\
$\forall x\in[0,1]$, $x\in[\frac{1}{n},\frac{2}{n}]$ for only finitely many $n$, hence $f_n(x)\to0$.\\
If $\alpha>0$ and $n\geq\alpha$, then $\lambda\{\abs{f_n}\geq\alpha\}=\frac{1}{n}\to0$, so $f_n\to0$ in measure.\\
Yet for $1\leq p<\infty$, $\norm{f_n}_p^p=n^p\lambda([\frac{1}{n},\frac{2}{n}])=n^{p-1}\nrightarrow0$.\\
Given $\delta>0$, choose $0<\eta<\delta$; on $[0,1]\setminus[0,\eta]$, $\forall n>\frac{2}{\eta}$, $f_n=0$, so $f_n\to0$ almost uniformly.\\
Thus a.e., in-measure, and almost-uniform convergence do not imply $L^p$ convergence.
\end{example}

\begin{theorem} \label{thm:pointwise-ae-norms-characterize-lp-convergence} \hlt{Pointwise A.E. Convergence Plus Norms}\\
Let $1\leq p<\infty$, $f_n,f\in L^p(\mu)$, $f_n\to f$ a.e. Then
$f_n\to f$ in $L^p(\mu)$ iff $\lim\limits_{n\to\infty}\int_X\abs{f_n}^p\,d\mu=\int_X\abs f^p\,d\mu$.
\begin{proof}
\begin{enumerate}[label=\roman*.]
\setlength{\itemsep}{0pt}
\item[$(\Rightarrow)$] If $\norm{f_n-f}_p\to0$, then $\abs{\norm{f_n}_p-\norm f_p}\leq\norm{f_n-f}_p$, hence $\int_X\abs{f_n}^p\,d\mu=\norm{f_n}_p^p\to\norm f_p^p=\int_X\abs f^p\,d\mu$.
\item[$(\Leftarrow)$] Assume $\int_X\abs{f_n}^p\,d\mu\to\int_X\abs f^p\,d\mu$. Set $h_n\defeq\frac{\abs{f_n}^p+\abs f^p}{2}-\abs{\frac{f_n-f}{2}}^p\geq0$ by convexity.\\
Since $h_n\to\abs f^p$ a.e., Fatou gives $\int_X\abs f^p\,d\mu\leq\liminf\limits_n\int_Xh_n\,d\mu=\int_X\abs f^p\,d\mu-2^{-p}\limsup\limits_n\int_X\abs{f_n-f}^p\,d\mu$.\\
Thus $\limsup\limits_n\norm{f_n-f}_p^p\leq0$, so $f_n\to f$ in $L^p(\mu)$.
\end{enumerate}
\hlt{Structural Commentary}: A.e. convergence identifies the limit; $p$-mass convergence prevents loss.\\
\hlt{Key Insight}: Fatou applied to the convexity defect forces the error norm to vanish.
\end{proof}
\end{theorem}

\begin{definition} \label{def:convergence-in-measure-and-cauchy} \hlt{Convergence and Cauchy Convergence in Measure}\\
Let $(f_n)$ be measurable functions on $X$.
\begin{enumerate}[label=(\roman*)]
\setlength{\itemsep}{0pt}
\item $f_n\to f$ \hlt{in measure} if $\forall\alpha>0$, $\mu\{x\in X:\abs{f_n(x)-f(x)}\geq\alpha\}\to0$.
\item $(f_n)$ is \hlt{Cauchy in measure} if $\forall\alpha>0$, $\mu\{x\in X:\abs{f_m(x)-f_n(x)}\geq\alpha\}\to0$ as $m,n\to\infty$.
\end{enumerate}
The same definitions result with $>$ in place of $\geq$.
\end{definition}

\begin{proposition} \label{prop:lp-convergence-implies-in-measure} \hlt{$L^p$ Convergence Implies Convergence in Measure}\\
Let $1\leq p<\infty$. If $\norm{f_n-f}_p\to0$, then $f_n\to f$ in measure.
\begin{proof}
For $\alpha>0$, set $E_n(\alpha)\defeq\{\abs{f_n-f}\geq\alpha\}$.\\
Since $\norm{f_n-f}_p^p\geq\int_{E_n(\alpha)}\abs{f_n-f}^p\,d\mu\geq\alpha^p\mu(E_n(\alpha))$, $\mu(E_n(\alpha))\leq\alpha^{-p}\norm{f_n-f}_p^p\to0$.\\
\hlt{Structural Commentary}: Large-error sets are controlled by the integral of the $p$th-power error.\\
\hlt{Key Insight}: This is Chebyshev's inequality in deterministic measure form.
\end{proof}
\end{proposition}

\begin{example} \label{ex:typewriter-lp-no-pointwise} \hlt{$L^p$ Convergence Without Pointwise Convergence}\\
Let $X=[0,1]$ with Lebesgue measure and list $I_{m,k}\defeq[\frac{k-1}{m},\frac{k}{m}]$, $1\leq k\leq m$, row by row in $m$.\\
Let $f_n\defeq\mathbf{1}_{I_n}$, where $I_n$ is the $n$th interval. If $I_n$ lies in row $m(n)$, then $m(n)\to\infty$ and\\
$\norm{f_n}_p^p=\lambda(I_n)=\frac{1}{m(n)}\to0$; hence $f_n\to0$ in $L^p$, therefore in measure by Proposition \ref{prop:lp-convergence-implies-in-measure}.\\
$\forall x\in[0,1]$, every row contains some $I_{m,k}$ with $x\in I_{m,k}$, giving a subsequence with value $1$ at $x$; for infinitely many rows, some $I_{m,k}$ omits $x$, giving a subsequence with value $0$ at $x$. Thus $(f_n(x))$ converges nowhere.\\
The subfamily $I_{m,2}=[\frac{1}{m},\frac{2}{m}]$, $m\geq2$, gives a subsequence converging pointwise to $0$ on all of $X$.
\end{example}

\begin{figure}[H]
  \centering
  % Source: attached Chapter 7, Example 7.4; clean redraw of the interval
% enumeration behind the typewriter sequence.
\begin{tikzpicture}[
  x=8.6cm,
  y=0.78cm,
  >=Stealth,
  every node/.style={font=\scriptsize},
  base/.style={black!65, line width=0.55pt},
  tick/.style={black!45, line width=0.45pt},
  interval/.style={fill=blue!8, draw=blue!45!black, line width=0.5pt},
  active/.style={fill=blue!28, draw=blue!70!black, line width=0.7pt},
  hit/.style={red!70!black, dashed, line width=0.7pt}
]
  \def\xhit{0.63}
  \foreach \m/\y in {1/4,2/3,3/2,4/1,5/0} {
    \draw[base] (0,\y) -- (1,\y);
    \node[left=7pt] at (0,\y) {$m=\m$};
    \pgfmathtruncatemacro{\hitj}{floor(\m*\xhit)+1}
    \foreach \j in {1,...,\m} {
      \pgfmathsetmacro{\a}{(\j-1)/\m}
      \pgfmathsetmacro{\b}{\j/\m}
      \ifnum\j=\hitj
        \filldraw[active] (\a,\y-0.13) rectangle (\b,\y+0.13);
      \else
        \filldraw[interval] (\a,\y-0.13) rectangle (\b,\y+0.13);
      \fi
      \draw[tick] (\a,\y-0.22) -- (\a,\y+0.22);
    }
    \draw[tick] (1,\y-0.22) -- (1,\y+0.22);
  }

  \draw[hit] (\xhit,-0.35) -- (\xhit,4.38);
  \node[red!70!black, above] at (\xhit,4.38) {fixed $x$};
  \node[below] at (0,-0.34) {$0$};
  \node[below] at (1,-0.34) {$1$};
\end{tikzpicture}

  \caption{In each row, one shrinking typewriter interval contains the fixed point \textcolor{red!70!black}{$x$}.}
  \label{fig:typewriter-intervals-lp-no-pointwise}
\end{figure}

\begin{theorem} \label{thm:riesz-subsequence-cauchy-measure} \hlt{Riesz Subsequence Principle for Convergence in Measure}\\
If $(f_n)$ is Cauchy in measure, then there exist a subsequence $(g_k)$ and a measurable $f$ such that:
\begin{enumerate}[label=(\roman*)]
\setlength{\itemsep}{0pt}
\item $g_k\to f$ a.e.
\item $g_k\to f$ in measure.
\item $\forall\varepsilon>0\ \exists E\in\mathcal{M}$ with $\mu(E)<\varepsilon$ such that $g_k\to f$ uniformly on $X\setminus E$.
\end{enumerate}
\begin{proof}
\begin{enumerate}[label=(\roman*)]
\setlength{\itemsep}{0pt}
\item Choose $n_1<n_2<\cdots$ so that $m,n\geq n_k\Rightarrow\mu\{\abs{f_m-f_n}\geq\frac{1}{2^k}\}<\frac{1}{2^k}$, and set $g_k\defeq f_{n_k}$.\\
For $E_k\defeq\{\abs{g_{k+1}-g_k}\geq\frac{1}{2^k}\}$, $\mu(E_k)<\frac{1}{2^k}$. Set $F_k\defeq\bigcup\limits_{j=k}^{\infty}E_j$. Then $F_k\downarrow$, $\mu(F_k)\leq\sum\limits_{j=k}^{\infty}\frac{1}{2^j}=\frac{1}{2^{k-1}}$.\\
For $i>j\geq k$, $x\notin F_k$, $\abs{g_i(x)-g_j(x)}\leq\sum\limits_{r=j}^{i-1}\abs{g_{r+1}(x)-g_r(x)}<\frac{1}{2^{j-1}}$.\\
Hence $(g_j)$ is uniformly Cauchy on $X\setminus F_k$ $\forall k$.\\
Let $F\defeq\bigcap\limits_kF_k$, so $\mu(F)=0$. Define $f=\lim\limits_jg_j$ on $X\setminus F$, $f=0$ on $F$. $f$ is measurable, $g_j\to f$ a.e.
\item Letting $i\to\infty$ gives $j\geq k$, $x\notin F_k\Rightarrow \abs{f(x)-g_j(x)}\leq\frac{1}{2^{j-1}}\leq\frac{1}{2^{k-1}}$. Given $\alpha,\varepsilon>0$, choose $k$ with $\frac{1}{2^{k-1}}<\alpha$ and $\mu(F_k)<\varepsilon$; then $j\geq k\Rightarrow \{\abs{f-g_j}\geq\alpha\}\subseteq F_k$, so $g_j\to f$ in measure.
\item Given $\varepsilon>0$, choose $k$ with $\mu(F_k)<\varepsilon$.\\
The estimate in (ii) gives $\sup\limits_{x\in X\setminus F_k}\abs{f(x)-g_j(x)}\leq\frac{1}{2^{j-1}}\to0$, so $g_j\to f$ uniformly on $X\setminus F_k$.
\end{enumerate}
\hlt{Structural Commentary}: Cauchy in measure lets us select summably small exceptional sets.\\
\hlt{Key Insight}: Outside the tail union of exceptional sets, the subsequence is uniformly Cauchy.
\end{proof}
\end{theorem}

\begin{corollary} \label{cor:cauchy-in-measure-completion} \hlt{Completeness for Convergence in Measure}\\
$(f_n)$ is Cauchy in measure iff $\exists f$ measurable with $f_n\to f$ in measure.
\begin{proof}
\begin{enumerate}[label=\roman*.]
\setlength{\itemsep}{0pt}
\item[$(\Rightarrow)$] If $(f_n)$ is Cauchy in measure, then $f_{n_k}\to f$ in measure for some measurable $f$ by Theorem \ref{thm:riesz-subsequence-cauchy-measure}.\\
For $\alpha>0$, $\{\abs{f_n-f}\geq\alpha\}\subseteq\{\abs{f_n-f_{n_k}}\geq\frac{\alpha}{2}\}\cup\{\abs{f_{n_k}-f}\geq\frac{\alpha}{2}\}$.\\
Choose $k$ large, then $n$ large, to get $f_n\to f$ in measure.
\item[$(\Leftarrow)$] For $\alpha>0$, $\{\abs{f_m-f_n}\geq\alpha\}\subseteq\{\abs{f_m-f}\geq\frac{\alpha}{2}\}\cup\{\abs{f_n-f}\geq\frac{\alpha}{2}\}$.\\
Hence $\mu\{\abs{f_m-f_n}\geq\alpha\}\to0$ as $m,n\to\infty$.
\end{enumerate}
\hlt{Structural Commentary}: The subsequential in-measure limit attracts the whole Cauchy sequence.\\
\hlt{Key Insight}: Convergence in measure is complete after measurable limits are admitted.
\end{proof}
\end{corollary}

\begin{corollary} \label{cor:cauchy-in-measure-completion-uniqueness} \hlt{Uniqueness for Convergence in Measure}\\
If $f_n\to f$ and $f_n\to g$ in measure, then $f=g$ a.e.
\begin{proof}
For $\alpha>0$, $\{\abs{f-g}\geq\alpha\}\subseteq\{\abs{f-f_n}\geq\frac{\alpha}{2}\}\cup\{\abs{f_n-g}\geq\frac{\alpha}{2}\}$, so $\mu\{\abs{f-g}\geq\alpha\}=0$.\\
Hence $\mu\{\abs{f-g}>0\}=0$ and $f=g$ a.e.\\
\hlt{Structural Commentary}: Two in-measure limits cannot differ on a positive-measure set.\\
\hlt{Key Insight}: In-measure limits are unique modulo null sets.
\end{proof}
\end{corollary}

\begin{corollary} \label{cor:in-measure-convergence-has-ae-subsequence} \hlt{In-Measure Limits Have A.E. Subsequences}\\
If $f_n\to f$ in measure, then some subsequence $f_{n_k}$ converges to $f$ a.e.
\begin{proof} Since $f_n\to f$ in measure, $(f_n)$ is Cauchy in measure.\\
Choose a subsequence $f_{n_k}\to h$ a.e. and in measure by Theorem \ref{thm:riesz-subsequence-cauchy-measure}.\\
Then $h=f$ a.e. by Corollary \ref{cor:cauchy-in-measure-completion-uniqueness}; after enlarging the exceptional null set, $f_{n_k}\to f$ a.e.\\
\hlt{Structural Commentary}: Convergence in measure contains an almost-everywhere convergent subsequence.\\
\hlt{Key Insight}: Select a Cauchy-in-measure subsequence, then identify its limit by uniqueness.
\end{proof}
\end{corollary}

\begin{proposition} \label{prop:subsequence-and-indicator-exercise-consequences} \hlt{Subsequence and Indicator Consequences}
\begin{enumerate}[label=(\roman*)]
\setlength{\itemsep}{0pt}
\item $f_n\to f$ in measure $\Rightarrow$ every subsequence $f_{n_k}\to f$ in measure; if $(f_n)$ is Cauchy in measure,\\
then every $(f_{n_k})$ is Cauchy in measure.
\item If $\norm{f_n-f}_p\to0$ and $\norm{f_{n_k}-g}_p\to0$, then $f=g$ a.e.
\item If $f_n=\mathbf{1}_{E_n}$ and $\norm{f_n-f}_p\to0$, then $f=\mathbf{1}_E$ a.e. for some $E\in\mathcal{M}$.
\end{enumerate}
\begin{proof}
\begin{enumerate}[label=(\roman*)]
\setlength{\itemsep}{0pt}
\item For $\alpha>0$, $\mu\{\abs{f_{n_k}-f}\geq\alpha\}\to0$ since $n_k\to\infty$; the Cauchy claim is identical with $m_k,n_\ell\to\infty$.
\item $f_n\to f$ and $f_{n_k}\to g$ in measure by Proposition \ref{prop:lp-convergence-implies-in-measure}; also $f_{n_k}\to f$ in measure by (i).\\
Thus $f=g$ a.e. by Corollary \ref{cor:cauchy-in-measure-completion-uniqueness}.
\item By Proposition \ref{prop:lp-convergence-implies-in-measure} and Theorem \ref{thm:riesz-subsequence-cauchy-measure}, some $f_{n_k}\to f$ a.e.\\
Since $f_{n_k}\in\{0,1\}$, $f\in\{0,1\}$ a.e.; with $E\defeq\{f=1\}$, $E\in\mathcal{M}$ and $f=\mathbf{1}_E$ a.e.
\end{enumerate}
\hlt{Structural Commentary}: Subsequence permanence and a.e. uniqueness transfer mode limits to representatives.\\
\hlt{Key Insight}: In-measure subsequences preserve limits; indicator limits remain indicators a.e.
\end{proof}
\end{proposition}

\begin{proposition} \label{prop:finite-measure-algebra-convergence-in-measure} \hlt{Algebra of Convergence in Measure on Finite Measure Spaces}\\
Let $(X,\mathcal{M},\mu)$ satisfy $\mu(X)<+\infty$. All functions below are measurable and finite-valued a.e.
\begin{enumerate}[label=(\roman*)]
\setlength{\itemsep}{0pt}
\item If $f_n\to f$ and $g_n\to g$ in measure, then $\forall a,b\in\R$, $af_n+bg_n\to af+bg$ in measure.
\item If $u_n\to0$ in measure and $h$ is measurable and finite-valued a.e., then $hu_n\to0$ in measure.
\item If $f_n\to f$ and $g_n\to g$ in measure, then $f_ng_n\to fg$ in measure.
\item If $f_n\to f$ in measure, then $f_n^2\to f^2$ in measure.
\end{enumerate}
\begin{proof}
\begin{enumerate}[label=(\roman*)]
\setlength{\itemsep}{0pt}
\item Fix $\alpha>0$. Since $\{\abs{a(f_n-f)+b(g_n-g)}\geq\alpha\}\subseteq\{\abs{a(f_n-f)}\geq\frac{\alpha}{2}\}\cup\{\abs{b(g_n-g)}\geq\frac{\alpha}{2}\}$, the measure of the right-hand side tends to $0$ by convergence in measure of $f_n$ and $g_n$.
\item Fix $\alpha,\eta>0$. Since $h$ is finite a.e. and $\mu(X)<+\infty$, choose $M>0$ with $\mu\{\abs h>M\}<\eta$.\\
Then $\{\abs{hu_n}\geq\alpha\}\subseteq\{\abs h>M\}\cup\{\abs{u_n}\geq\frac{\alpha}{M}\}$, so $\limsup\limits_n\mu\{\abs{hu_n}\geq\alpha\}\leq\eta$. Let $\eta\downarrow0$.
\item Set $u_n\defeq f_n-f$ and $v_n\defeq g_n-g$. Then $u_n,v_n\to0$ in measure and $f_ng_n-fg=fv_n+gu_n+u_nv_n$.\\
By (ii), $fv_n\to0$ and $gu_n\to0$ in measure. For $\alpha>0$, $\{\abs{u_nv_n}\geq\alpha\}\subseteq\{\abs{u_n}\geq\sqrt{\alpha}\}\cup\{\abs{v_n}\geq\sqrt{\alpha}\}$, so $u_nv_n\to0$ in measure. Apply (i) to the three summands.
\item Apply (iii) with $g_n=f_n$ and $g=f$.
\end{enumerate}
\hlt{Structural Commentary}: Finite measure makes finite measurable multipliers harmless after truncation.\\
\hlt{Key Insight}: Products split into fixed-multiplier errors and a product of two small-in-measure errors.
\end{proof}
\end{proposition}

\begin{theorem} \label{thm:dominated-measure-to-lp} \hlt{Dominated Convergence in Measure Gives $L^p$ Convergence}\\
Let $1\leq p<\infty$, $f_n\in L^p(\mu)$, $f_n\to f$ in measure, and $\abs{f_n}\leq g$ a.e. for some $g\in L^p(\mu)$. Then:
\begin{enumerate}[label=(\roman*)]
\setlength{\itemsep}{0pt}
\item $f\in L^p(\mu)$;
\item $\norm{f_n-f}_p\to0$.
\end{enumerate}
\begin{proof}
Since $f_n\to f$ in measure, $(f_n)$ is Cauchy in measure.\\
\begin{enumerate}[label=(\roman*)]
\setlength{\itemsep}{0pt}
\item By Theorem \ref{thm:riesz-subsequence-cauchy-measure}, some subsequence $(u_j)$ satisfies $u_j\to h$ a.e. and in measure; $h=f$ a.e. by Corollary \ref{cor:cauchy-in-measure-completion-uniqueness}. Since $\abs{u_j}\leq g$ a.e., $f\in L^p(\mu)$ by Theorem \ref{thm:dominated-ae-to-lp}.
\item If norm convergence fails, then $\exists\varepsilon>0$ and a subsequence $(g_k)$ with $\norm{g_k-f}_p\geq\varepsilon$. The argument in (i) produces a subsubsequence $(h_j)$ with $h_j\to f$ a.e.; then $\norm{h_j-f}_p\to0$ by Theorem \ref{thm:dominated-ae-to-lp}, a contradiction.
\end{enumerate}
\hlt{Structural Commentary}: Reduce convergence in measure to almost-everywhere convergence along a subsequence.\\
\hlt{Key Insight}: Domination turns every in-measure subsequence into an $L^p$-convergent subsubsequence.
\end{proof}
\end{theorem}

\begin{lemma} \label{lem:almost-uniform-cauchy-has-ae-limit} \hlt{Almost Uniform Cauchy Sequences Have Almost Uniform Limits}\\
If $(f_n)$ is almost uniformly Cauchy, then $\exists f$ measurable such that $f_n\to f$ almost uniformly and a.e.
\begin{proof}
Choose $E_k\in\mathcal{M}$ with $\mu(E_k)<2^{-k}$ and $(f_n)$ uniformly Cauchy on $X\setminus E_k$. Set $F_k\defeq\bigcup\limits_{j=k}^{\infty}E_j$ and $F\defeq\bigcap\limits_kF_k$; then $F_k\downarrow$, $\mu(F_k)<2^{-(k-1)}$, $\mu(F)=0$, and $(f_n)$ is uniformly Cauchy on $X\setminus F_k$. By Theorem \ref{thm:uniform-cauchy-criterion}, $f_n$ has a uniform limit on each $X\setminus F_k$; these limits agree on overlaps, so define measurable $f$ on $X\setminus F$, and set $f=0$ on $F$. Hence $f_n\to f$ a.e. and uniformly on each $X\setminus F_k$; since $\mu(F_k)\to0$, $f_n\to f$ almost uniformly.\\
\hlt{Structural Commentary}: Tail unions convert many small exceptional sets into nested uniform domains.\\
\hlt{Key Insight}: Almost uniform Cauchy means uniformly Cauchy after discarding a tail-small exceptional set.
\end{proof}
\end{lemma}

\begin{theorem} \label{thm:almost-uniform-and-measure-subsequences} \hlt{Almost Uniform Convergence and Convergence in Measure}\\
The following hold.
\begin{enumerate}[label=(\roman*)]
\setlength{\itemsep}{0pt}
\item If $f_n\to f$ almost uniformly, then $f_n\to f$ in measure.
\item If $f_n\to f$ in measure, then $\exists(f_{n_k})$ with $f_{n_k}\to f$ almost uniformly.
\end{enumerate}
\begin{proof}
\begin{enumerate}[label=(\roman*)]
\setlength{\itemsep}{0pt}
\item If $f_n\rightrightarrows f$ on $X\setminus E_\delta$ with $\mu(E_\delta)<\delta$, then $\forall\alpha>0$ and sufficiently large $n$, $\{\abs{f_n-f}\geq\alpha\}\subseteq E_\delta$.\\
Thus $\limsup\limits_n\mu\{\abs{f_n-f}\geq\alpha\}\leq\delta$, and let $\delta\downarrow0$.
\item Convergence in measure makes $(f_n)$ Cauchy in measure. Choose $f_{n_k}\to h$ almost uniformly and in measure by Theorem \ref{thm:riesz-subsequence-cauchy-measure}. Uniqueness gives $h=f$ a.e.; absorb $\{h\neq f\}$ into the exceptional sets.
\end{enumerate}
\hlt{Structural Commentary}: Almost uniform convergence is uniform convergence off arbitrarily small sets.\\
\hlt{Key Insight}: Convergence in measure always hides an almost-uniform subsequence.
\end{proof}
\end{theorem}

\begin{corollary} \label{cor:finite-measure-ae-implies-measure} \hlt{Finite-Measure Almost-Everywhere Convergence Implies In-Measure Convergence}\\
If $\mu(X)<+\infty$ and $f_n\to f$ a.e., with $f_n,f$ finite-valued off a null set, then $f_n\to f$ in measure.
\begin{proof}
By Theorem \ref{thm:egoroff-finite-measure}, $f_n\to f$ almost uniformly; hence $f_n\to f$ in measure by Theorem \ref{thm:almost-uniform-and-measure-subsequences}.\\
\hlt{Structural Commentary}: Finite measure upgrades a.e. convergence through almost uniform convergence.\\
\hlt{Key Insight}: On finite measure spaces, a.e. convergence implies in-measure convergence.
\end{proof}
\end{corollary}

\begin{corollary} \label{cor:finite-measure-subsequence-test-convergence-in-measure} \hlt{Finite-Measure Subsequence Test for Convergence in Measure}\\
Let $(X,\mathcal{M},\mu)$ satisfy $\mu(X)<+\infty$, and let $f_n,f:X\to\overline{\R}$ be measurable and finite-valued a.e.\\
Then $f_n\to f$ in measure iff every subsequence of $(f_n)$ has a further subsequence converging to $f$ a.e.
\begin{proof}
\begin{enumerate}[label=\roman*.]
\setlength{\itemsep}{0pt}
\item[$(\Rightarrow)$] If $f_n\to f$ in measure, then every subsequence also converges to $f$ in measure by Proposition \ref{prop:subsequence-and-indicator-exercise-consequences}.\\
Choose a further subsequence converging to $f$ a.e. by Corollary \ref{cor:in-measure-convergence-has-ae-subsequence}.
\item[$(\Leftarrow)$] Suppose $f_n\nrightarrow f$ in measure. Then $\exists \alpha,\varepsilon>0$ and subsequence $(f_{n_k})$ with $\mu\{\abs{f_{n_k}-f}\geq\alpha\}\geq\varepsilon$ $\forall k$.\\
But some $f_{n_{k_j}}\to f$ a.e., hence in measure by Corollary \ref{cor:finite-measure-ae-implies-measure}, contradicting the uniform lower bound.
\end{enumerate}
\hlt{Structural Commentary}: Subsequence tests characterize in-measure convergence.\\
\hlt{Key Insight}: Nonconvergence leaves a bad subsequence.
\end{proof}
\end{corollary}

\begin{figure}[H]
  \centering
  % Source: attached Chapter 7, Diagrams 7.1--7.3; simplified redraw.
% Solid arrows are implications. Dashed arrows mean a subsequence converges in
% the indicated mode.
\begin{tikzpicture}[
  scale=0.94,
  >=Stealth,
  every node/.style={font=\scriptsize},
  mode/.style={draw=blue!60!black, fill=blue!6, rounded corners=2pt,
    text=black,
    minimum width=0.72cm, minimum height=0.42cm, inner sep=2pt},
  title/.style={font=\scriptsize\itshape, black!75},
  imp/.style={->, black!78, line width=0.8pt},
  eqimp/.style={<->, black!78, line width=0.8pt},
  sub/.style={->, red!65!black, dashed, line width=0.75pt}
]
  \begin{scope}[xshift=0cm]
    \node[title] at (0,1.50) {general};
    \node[mode] (gAE) at (-1.05,0.62) {$AE$};
    \node[mode] (gAU) at (1.05,0.62) {$AU$};
    \node[mode] (gLP) at (-1.05,-0.62) {$L^p$};
    \node[mode] (gM) at (1.05,-0.62) {$M$};

    \draw[imp] (gAU) -- (gAE);
    \draw[imp] (gAU) -- (gM);
    \draw[imp] (gLP) -- (gM);
    \draw[sub] (gM) to[bend left=18] (gAU);
  \end{scope}

  \begin{scope}[xshift=3.65cm]
    \node[title] at (0,1.50) {finite measure};
    \node[mode] (fAE) at (-1.05,0.62) {$AE$};
    \node[mode] (fAU) at (1.05,0.62) {$AU$};
    \node[mode] (fLP) at (-1.05,-0.62) {$L^p$};
    \node[mode] (fM) at (1.05,-0.62) {$M$};

    \draw[eqimp] (fAE) -- (fAU);
    \draw[imp] (fAU) -- (fM);
    \draw[imp] (fLP) -- (fM);
    \draw[sub] (fM) to[bend left=18] (fAU);
  \end{scope}

  \begin{scope}[xshift=7.30cm]
    \node[title] at (0,1.50) {dominated by $g\in L^p$};
    \node[mode] (dAE) at (-1.05,0.62) {$AE$};
    \node[mode] (dAU) at (1.05,0.62) {$AU$};
    \node[mode] (dLP) at (-1.05,-0.62) {$L^p$};
    \node[mode] (dM) at (1.05,-0.62) {$M$};

    \draw[imp] (dAU) -- (dAE);
    \draw[imp] (dAE) -- (dLP);
    \draw[eqimp] (dLP) -- (dM);
    \draw[sub] (dM) to[bend left=18] (dAU);
  \end{scope}
\end{tikzpicture}

  \caption{Modes of convergence in three settings: solid arrows are implications; dashed arrows yield convergent subsequences.}
  \label{fig:modes-convergence-implication-diagrams}
\end{figure}

\begin{method} \label{meth:modes-of-convergence-counterexample-dictionary} \hlt{Counterexample Dictionary}
\begin{enumerate}[label=(\roman*)]
\setlength{\itemsep}{0pt}
\item On $(\R,\lambda)$, $f_n=n^{-\frac{1}{p}}\mathbf{1}_{[0,n]}$ satisfies $f_n\rightrightarrows0$ but $\norm{f_n}_p=1$.
\item On $([0,1],\lambda)$, $f_n=n\mathbf{1}_{[\frac{1}{n},\frac{2}{n}]}$ satisfies $f_n\to0$ a.e., in measure, almost uniformly, but $\norm{f_n}_p^p=n^{p-1}\nrightarrow0$.
\item Typewriter intervals on $[0,1]$: if $f_n=\mathbf{1}_{I_n}$ and $\lambda(I_n)\to0$ rowwise, then $\norm{f_n}_p\to0$ and $f_n\to0$ in measure, but $(f_n(x))$ converges nowhere.
\item On $(\R,\lambda)$, $f_n=\mathbf{1}_{[n,n+1]}$ satisfies $f_n\to0$ pointwise but $\forall n,\ \lambda\{\abs{f_n}\geq\frac{1}{2}\}=1$, so $f_n\nrightarrow0$ in measure.
\item On $(\R,\lambda)$, $f_n=\mathbf{1}_{[n,n+1]}$ is not almost uniformly convergent to $0$: if $\lambda(E)<1$ and $f_n\rightrightarrows0$ on $\R\setminus E$, then $[n,n+1]\subseteq E$ for sufficiently large $n$, impossible.
\end{enumerate}
\end{method}

\begin{proposition} \label{prop:finite-measure-integral-metric-for-measure-convergence} \hlt{An Integral Metric for Convergence in Measure on Finite Measure Spaces}\\
Let $\mu(X)<+\infty$ and $r(h)\defeq\int_X\abs{h}(1+\abs{h})^{-1}\,d\mu$. Then $f_n\to f$ in measure iff $r(f_n-f)\to0$.
\begin{proof}
\begin{enumerate}[label=\roman*.]
\setlength{\itemsep}{0pt}
\item[$(\Rightarrow)$] Set $h_n\defeq f_n-f$. For $\eta>0$ and $A_n\defeq\{\abs{h_n}\geq\eta\}$, $r(h_n)\leq\mu(A_n)+\eta\mu(X)$. Let $n\to\infty$, then $\eta\downarrow0$.
\item[$(\Leftarrow)$] For $\alpha>0$, $\alpha(1+\alpha)^{-1}\mathbf{1}_{\{\abs{h_n}\geq\alpha\}}\leq\abs{h_n}(1+\abs{h_n})^{-1}$, hence $\mu\{\abs{h_n}\geq\alpha\}\leq(1+\alpha)\alpha^{-1}r(h_n)\to0$.
\end{enumerate}
\hlt{Structural Commentary}: The bounded error transform metrizes convergence in measure on finite measure spaces.\\
\hlt{Key Insight}: Small integral of the bounded error transform is equivalent to small large-error sets.
\end{proof}
\end{proposition}

\begin{corollary} \label{cor:fatou-and-dct-with-convergence-in-measure} \hlt{Fatou and Dominated Convergence with In-Measure Hypotheses}
\begin{enumerate}[label=(\roman*)]
\setlength{\itemsep}{0pt}
\item $f_n\geq0$ and $f_n\to f$ in measure $\Rightarrow \int_X f\,d\mu\leq\liminf\limits_n\int_Xf_n\,d\mu$.
\item $f_n\to f$ in measure and $\abs{f_n}\leq g\in\mathscr{L}^1(\mu)$ $\Rightarrow \int_Xf_n\,d\mu\to\int_Xf\,d\mu$.
\end{enumerate}
\begin{proof}
\begin{enumerate}[label=(\roman*)]
\setlength{\itemsep}{0pt}
\item Choose $f_{n_k}$ with $\int f_{n_k}\,d\mu\to\liminf\limits_n\int f_n\,d\mu$. By Theorem \ref{thm:riesz-subsequence-cauchy-measure} and uniqueness, some $f_{n_{k_j}}\to f$ a.e.\\
Fatou gives $\int f\,d\mu\leq\liminf\limits_j\int f_{n_{k_j}}\,d\mu=\liminf\limits_n\int f_n\,d\mu$.
\item By Theorem \ref{thm:dominated-measure-to-lp} with $p=1$, $\norm{f_n-f}_1\to0$, hence $\abs{\int_X(f_n-f)\,d\mu}\leq\norm{f_n-f}_1\to0$.
\end{enumerate}
\hlt{Structural Commentary}: In-measure hypotheses reduce to a.e. subsequences or $L^1$ convergence.\\
\hlt{Key Insight}: Fatou and DCT survive convergence in measure through subsequences and domination.
\end{proof}
\end{corollary}

The next criterion packages the local mass controls needed to reverse convergence in measure into norm convergence.

\begin{theorem} \label{thm:vitali-convergence-lp} \hlt{Vitali Convergence Criterion for $L^p$}\\
Let $1\leq p<\infty$, $f_n\in L^p(\mu)$, and let $f$ be measurable. Consider the following conditions.
\begin{enumerate}[label=(\roman*)]
\setlength{\itemsep}{0pt}
\item $f\in L^p(\mu)$ and $\norm{f_n-f}_p\to0$.
\item $f_n\to f$ in measure.
\item $\forall\varepsilon>0$ $\exists E_\varepsilon\in\mathcal{M}$ with $\mu(E_\varepsilon)<+\infty$ such that $F\cap E_\varepsilon=\emptyset\Rightarrow\sup\limits_n\int_F\abs{f_n}^p\,d\mu<\varepsilon^p$.
\item $\forall\varepsilon>0$ $\exists\delta(\varepsilon)>0$ such that $\mu(E)<\delta(\varepsilon)\Rightarrow\sup\limits_n\int_E\abs{f_n}^p\,d\mu<\varepsilon^p$.
\end{enumerate}
Precisely, (i) is equivalent to the conjunction of (ii)--(iv).
\begin{proof}
\begin{enumerate}[label=(\roman*)]
\setlength{\itemsep}{0pt}
\item Assume (i). Then (ii) holds by Proposition \ref{prop:lp-convergence-implies-in-measure}, and $(f_n)$ is $L^p$-Cauchy. Conditions (iii)--(iv) follow from Proposition \ref{prop:lp-cauchy-local-mass-control} after raising its norm bounds to the $p$th power.
\item Assume (ii)--(iv). Fix $\eta>0$, choose $E_\eta$ from (iii) and $\delta(\eta)$ from (iv), and put $F=X\setminus E_\eta$.\\
Then $\norm{(f_n-f_m)\mathbf{1}_F}_p\leq\norm{f_n\mathbf{1}_F}_p+\norm{f_m\mathbf{1}_F}_p<2\eta$.\\
If $\mu(E_\eta)>0$, set $\alpha\defeq\eta\mu(E_\eta)^{-\frac{1}{p}}$ and $H_{nm}\defeq E_\eta\cap\{\abs{f_n-f_m}\geq\alpha\}$.\\
By (ii), $(f_n)$ is Cauchy in measure, so $\mu(H_{nm})<\delta(\eta)$ for large $m,n$. On $E_\eta\setminus H_{nm}$, $\abs{f_n-f_m}<\alpha$; on $H_{nm}$, use $\abs{f_n-f_m}\leq\abs{f_n}+\abs{f_m}$. Hence $\norm{(f_n-f_m)\mathbf{1}_{E_\eta}}_p\leq\alpha\mu(E_\eta)^{\frac{1}{p}}+\norm{f_n\mathbf{1}_{H_{nm}}}_p+\norm{f_m\mathbf{1}_{H_{nm}}}_p<3\eta$ by (iv). If $\mu(E_\eta)=0$, the $E_\eta$ term is $0$. Thus $\norm{f_n-f_m}_p<5\eta$ eventually, so $(f_n)$ is $L^p$-Cauchy.\\
By Riesz--Fischer, $f_n\to h$ in $L^p$; by Proposition \ref{prop:lp-convergence-implies-in-measure} and Corollary \ref{cor:cauchy-in-measure-completion-uniqueness}, $h=f$ a.e., so (i) holds.
\end{enumerate}
\hlt{Structural Commentary}: Vitali = in-measure convergence + tightness + uniform small-set control.\\
\hlt{Key Insight}: Tightness and small-set control convert in-measure convergence into an $L^p$ Cauchy estimate.
\end{proof}
\end{theorem}
